%! TeX program = lualatex
\documentclass[
	% -- opções da classe memoir --
	12pt,				% tamanho da fonte
	oneside,			% para impressão em recto e verso. Oposto a oneside
	a4paper,			% tamanho do papel. 
	% -- opções da classe abntex2 --
	chapter=TITLE,		% títulos de capítulos convertidos em letras maiúsculas
	section=TITLE,		% títulos de seções convertidos em letras maiúsculas
	subsection=TITLE,	% títulos de subseções convertidos em letras maiúsculas
	subsubsection=TITLE,% títulos de subsubseções convertidos em letras maiúsculas
	% -- opções do pacote babel --
	english,			% idioma adicional para hifenização
	french,				% idioma adicional para hifenização
	spanish,			% idioma adicional para hifenização
	brazilian,				% o último idioma é o principal do documento
	]{abntex2}

% ---
% PACOTES
% ---

% ---
% Pacotes fundamentais 
% ---
\usepackage[T1]{fontenc}		% Selecao de codigos de fonte.
\usepackage[utf8]{inputenc}		% Codificacao do documento (conversão automática dos acentos)
\usepackage{indentfirst}		% Indenta o primeiro parágrafo de cada seção.
\usepackage{color}				% Controle das cores
\usepackage{graphicx}			% Inclusão de gráficos
\usepackage{microtype} 			% para melhorias de justificação
\usepackage{fontspec}
\setmainfont{Arial}
\setsansfont{Arial}

% Regras da ABNT — títulos em 12pt (normalsize deste documento), negrito
\renewcommand{\ABNTEXchapterfont}{\sffamily\bfseries}
\renewcommand{\ABNTEXpartfont}{\ABNTEXchapterfont}
\renewcommand{\ABNTEXchapterfontsize}{\normalsize}
\renewcommand{\ABNTEXsectionfontsize}{\normalsize}
\renewcommand{\ABNTEXsubsectionfontsize}{\normalsize}
\renewcommand{\ABNTEXsubsubsectionfontsize}{\normalsize}

% ---
% Cabecalhos: apenas o número da página, sem o risco separador
% (baseado no estilo do abntex2, removendo o marcador de capítulo/seção)
% ---
\makeevenhead{abntheadings}{\ABNTEXfontereduzida\thepage}{}{}
\makeoddhead{abntheadings}{}{}{\ABNTEXfontereduzida\thepage}
\makeheadrule{abntheadings}{0pt}{0pt} % remove o risco abaixo do cabeçalho

% ---
% Pacotes de citações
% ---
\usepackage[brazilian,hyperpageref]{backref}	 % Paginas com as citações na bibl
\usepackage[alf]{abntex2cite}


% ---
% Configurações do pacote backref
% Usado sem a opção hyperpageref de backref
\renewcommand{\backrefpagesname}{}
% Texto padrão antes do número das páginas
\renewcommand{\backref}{}
% Define os textos da citação
\renewcommand*{\backrefalt}[4]{}
% ---

% ---
% Informações de dados para CAPA e FOLHA DE ROSTO
% ---
\titulo{Usabilidade e carga de trabalho no registro de dados ecológicos: desenvolvimento de um sistema para triagem taxonômica de comunidades marinhas e estuarinas}
\autor{Gabriel Kavata Leandro}
\orientador{Guilherme Corredato Guerino}
\coorientador{Pablo Damian Borges Guilherme}
\local{Brasil}
\data{2026}
\instituicao{
  Universidade Estadual do Paraná -- UNESPAR
  \par
  Programa de Pós-Graduação em Ambientes Litorâneos e Insulares}
\tipotrabalho{Plano de trabalho (Mestrado)}
\preambulo{Plano de trabalho apresentado ao Programa de Pós-Graduação em Ambientes Litorâneos e Insulares, como requisito de ingresso ao curso de mestrado.}

% Alterando o aspecto da cor azul
\definecolor{blue}{RGB}{41,5,195}

% informações do PDF
\makeatletter
\hypersetup{
     	%pagebackref=true,
		pdftitle={\@title}, 
		pdfauthor={\@author},
    	pdfsubject={\imprimirpreambulo},
	    pdfcreator={LaTeX with abnTeX2},
		pdfkeywords={plano de trabalho}{ciência ambiental}{desenvolvimento de sistemas}, 
		colorlinks=true,       		% false: boxed links; true: colored links
    	linkcolor=blue,          	% color of internal links
    	citecolor=blue,        		% color of links to bibliography
    	filecolor=magenta,      		% color of file links
		urlcolor=blue,
		bookmarksdepth=4
}
\makeatother
% --- 
% ---
% Possibilita criação de Quadros e Lista de quadros.
% Ver https://github.com/abntex/abntex2/issues/176
%
\newcommand{\quadroname}{Quadro}
\newcommand{\listofquadrosname}{Lista de quadros}

\newfloat[chapter]{quadro}{loq}{\quadroname}
\newlistof{listofquadros}{loq}{\listofquadrosname}
\newlistentry{quadro}{loq}{0}

% configurações para atender às regras da ABNT
\setfloatadjustment{quadro}{\centering}
\counterwithout{quadro}{chapter}
\renewcommand{\cftquadroname}{\quadroname\space} 
\renewcommand*{\cftquadroaftersnum}{\hfill--\hfill}

\setfloatlocations{quadro}{hbtp} % Ver https://github.com/abntex/abntex2/issues/176

% --- 
% Espaçamentos entre linhas e parágrafos 
% --- 

% O tamanho do parágrafo é dado por:
\setlength{\parindent}{1.3cm}

% Controle do espaçamento entre um parágrafo e outro:
\setlength{\parskip}{0.2cm}  % tente também \onelineskip

% Espaçamento após o título de capítulo igual ao de seção (parskip)
\setlength{\afterchapskip}{\parskip}
\setlength{\beforechapskip}{\parskip}
% ---

% Início do documento
% ----
\begin{document}
\selectlanguage{brazilian}
\frenchspacing 

% ----------------------------------------------------------
% ----------------------------------------------------------
% ELEMENTOS TEXTUAIS
% ----------------------------------------------------------

\textual

\chapter{PLANO DE TRABALHO}

\noindent \textbf{Título da pesquisa}: \imprimirtitulo

\noindent \textbf{Nome do Candidato}: \imprimirautor

\noindent \textbf{Nome do Orientador 1}: \imprimirorientador

\noindent \textbf{Nome do Orientador 2}: \imprimircoorientador

\noindent \textbf{Linha de Pesquisa pretendida}: Desenvolvimento socioambiental e tecnológico em ambientes litorâneos e insulares

\section{INTRODUÇÃO}

As comunidades ecológicas são conjuntos de todas as populações de diferentes espécies de plantas, animais e microrganismos vivendo e interagindo em uma área \cite{miller_ciencia_2021}. 
O estudo de tais comunidades é um elemento central da ecologia, essencial para compreender a formação, a manutenção e as consequências da diversidade biológica dos ecossistemas \cite{morin_community_2011}. 
Essa compreensão fornece a base para entender como os processos ecológicos influenciam diretamente as populações humanas \cite{morin_community_2011}.

Entre os ecossistemas mais dinâmicos e suscetíveis a essas transformações estão os estuários. Por funcionarem como zonas de transição entre ambientes continentais e marinhos, o estudo das comunidades nesses locais é crucial para mensurar o impacto das ações antrópicas sobre a biodiversidade \cite{whitfield_fishes_2002}. 
Um exemplo dessa dinâmica é o Complexo Estuarino de Paranaguá (CEP), localizado no litoral do Paraná e composto pelas baías de Paranaguá, Antonina e Laranjeiras \cite{angeli_arsenic_2013}. 
Esse sistema é cenário de diferentes atividades antrópicas -- como a pesca, o turismo e a indústria portuária -- ao mesmo tempo em que abriga regiões de elevada riqueza biológica que encontram-se em situação de vulnerabilidade \cite{miura_mapping_2020}.

Nesse contexto, monitorar as comunidades ecológicas do CEP permite acompanhar a saúde desse ecossistema e garantir a sustentabilidade das atividades antrópicas que impulsionam a economia e a cultura local \cite{andriguetto_filho_diagnostico_2006,miura_mapping_2020}.
Estudos prévios evidenciam a necessidade desse acompanhamento: \citeonline{angeli_arsenic_2013} detectaram a presença de metais pesados em peixes da região, enquanto \citeonline{liebezeit_ddt_2011} registraram concentrações consideráveis do inseticida diclorodifeniltricloroetano (DDT) em peixes, esponjas e bivalves.

Apesar da importância desse monitoramento, observa-se que os métodos empregados na triagem, catalogação e análise taxonômica em comunidades ecológicas ainda são eminentemente manuais e carecem de gestão de dados \cite{angeli_arsenic_2013,liebezeit_ddt_2011,santos_diversidade_2024}.
Esse gargalo foi corroborado em entrevistas com um grupo de pesquisa em comunidades bentônicas da Universidade Estadual do Paraná (UNESPAR), no qual pesquisadores relataram dificuldades no registro e no processamento de informações ecológicas. Os relatos destacam desde problemas com a inclusão de táxons identificados em planilhas até a atualização e extração de dados das mesmas. A falta de um sistema estruturado compromete a padronização das informações, elevando o tempo de análise e prejudicando a qualidade dos resultados \cite{specht_data_2015}.

Esse problema torna-se mais crítico à medida que a ecologia passa a demandar volumes crescentes de dados, estruturados e reutilizáveis, para sustentar análises comparativas e de longo prazo.
Segundo \citeonline{michener_ecoinformatics_2012}, a ecologia está gradativamente se tornando uma ciência intensiva em dados. Isso decorre do fato de que tal ciência passa a se fundamentar em grandes volumes de informações, oriundas de fontes diversas, como a observação humana ou a captura por sensores \cite{kelling_data-intensive_2009}.
Nessa mesma direção, \citeonline{strasser_promoting_2011} argumenta que o processo de uma pesquisa ecológica pode ser compreendido como um ``ciclo de vida'' de dados, o qual compreende desde a coleta e organização dos dados até a sua descrição e depósito em um sistema repositório, englobando a preservação de longo prazo e culminando na descoberta, integração e análise dessas informações.

Frente a esse cenário, a Interação Humano-Computador (IHC) oferece o referencial adequado para orientar o desenvolvimento de um sistema que se ajuste ao fluxo de trabalho já existente do pesquisador, em vez de impor um modelo genérico de gestão de dados alheio a essa prática \cite{silva_interacao_2010}.
Dentro desse referencial, o \textit{Design Thinking} (DT) constitui um processo iterativo e centrado no usuário, estruturado nas fases de empatia, definição, ideação, prototipação e teste, voltado à concepção de soluções ajustadas a problemas complexos e mal definidos \cite{brown_design_2021}. Essa condição caracteriza justamente o cenário descrito por \citeonline{marinoni_colecoes_2025}, em que a ausência de sistemas de gestão não decorre da inexistência de soluções de software, mas de sua inadequação ao trabalho real do pesquisador.

Para que a solução projetada se ajuste a esse trabalho real, utilizam-se duas técnicas complementares e sequenciais de representação do usuário na estrutura do DT \cite{siricharoen_using_2021}: o mapa de empatia, produzido na fase de empatia, organiza os dados da pesquisa em torno do que o usuário pensa, sente, vê, ouve e faz, permitindo compreender sua experiência e suas intenções para além do registro de comportamentos; e a persona, que na fase de definição consolida essas evidências em um perfil arquetípico construído a partir de dados empíricos, tratando o usuário como uma personagem coerente -- e não como um conjunto de especificações --, o que facilita a comunicação das decisões de design. Dessas evidências, a definição converge para o Point of View (POV): uma formulação estruturada do problema que articula a descrição do usuário, a necessidade identificada, a razão subjacente e o critério de sucesso -- muitas vezes não explicitada -- que torna essa necessidade relevante \cite{rodriguez_point_2019}. O POV evita a generalidade de um problema mal definido e a especificidade prematura de uma solução, delimitando o espaço de ideação de soluções \cite{rodriguez_point_2019}.

A análise temática constitui um método para identificar, analisar e relatar padrões -- ou temas -- dentro de um conjunto de dados qualitativos, sendo suficientemente flexível para ser aplicado a diferentes tamanhos de amostra, referenciais teóricos e tipos de pergunta de pesquisa \cite{braun_using_2006}. Tal método organiza-se em fases recursivas, e não estritamente lineares: familiarização com os dados, geração de códigos iniciais, busca por temas a partir do agrupamento de códigos relacionados, revisão dos temas quanto à sua coerência interna e à distinção em relação aos demais, definição e nomeação de cada tema, e, por fim, produção do relatório analítico \cite{braun_using_2006}. Nesse processo, os temas emergem tanto do conteúdo das entrevistas quanto das questões que orientam a pesquisa, fornecendo a base empírica para as etapas subsequentes de definição do usuário e do problema.

Ao converter o problema delimitado em decisões de engenharia, define-se a especificação de requisitos do sistema. Segundo \citeonline{sommerville_software_2016}, os requisitos funcionais especificam os serviços que o sistema deve oferecer e seu comportamento frente a eventos externos -- aqui, a gestão dos dados de comunidades ecológicas. Os requisitos não funcionais, por sua vez, estabelecem as restrições e os atributos de qualidade que a solução deve satisfazer -- desempenho, usabilidade, confiabilidade e manutenibilidade --, sendo tão determinantes quanto os funcionais para que o sistema se ajuste ao fluxo de trabalho do pesquisador \cite{sommerville_software_2016}. O produto dessa atividade é o documento de especificação de requisitos, referência para o desenvolvimento do sistema.

A avaliação de tais soluções, por sua vez, apoia-se nos construtos de eficácia e eficiência definidos pela norma ISO 9241-11 \cite{iso9241_11_2018}, complementados pela mensuração da carga de trabalho mental exigida do usuário por meio do \textit{NASA Task Load Index} (NASA-TLX), instrumento que avalia a demanda percebida em seis subescalas: demanda mental, demanda física, demanda temporal, desempenho, esforço e frustração \cite{hart_development_1988}. A aplicação conjunta desses instrumentos permite verificar se um sistema cumpre a tarefa a que se propõe e, também, se o faz reduzindo o esforço cognitivo exigido do pesquisador em relação ao método atualmente empregado.

Diante do exposto, justifica-se o desenvolvimento de um sistema de registro orientado pelos princípios da IHC e estruturado segundo o processo de DT. O sistema proposto visa atuar nas etapas de coleta, descrição e preservação do ciclo de vida de dados, tratando diretamente o gargalo identificado na análise e levantamento de resultados científicos a partir do estudo de comunidades. Essa proposta tem como expectativa a redução do tempo e do esforço cognitivo despendidos pelo pesquisador, sem comprometer -- e possivelmente aumentando -- a eficácia desse processo.


\section{OBJETIVO}
O trabalho proposto tem por objetivo geral o desenvolvimento de um sistema para registro de comunidades ecológicas marinhas e estuarinas, visando investigar a eficácia do sistema em agilizar o processo de coleta e pesquisa dessas comunidades. Dessa forma, os objetivos específicos são:
\begin{itemize}
\item Realizar entrevistas com pesquisadores de comunidades relacionadas ao trabalho proposto;
\item Levantar requisitos funcionais e não funcionais para desenvolvimento do sistema;
\item Coletar dados de eficiência dos métodos de triagem e identificação de táxons antes do uso do sistema;
\item Desenvolver um sistema para registro de comunidades ecológicas marinhas e estuarinas;
\item Coletar dados de eficiência dos métodos de triagem e identificação de táxons durante o uso do sistema;
\item Comparar os dados coletados.
\end{itemize}

\section{METODOLOGIA}
Quanto à natureza, esta pesquisa é aplicada, pois visa gerar conhecimento dirigido à solução de um problema prático: agilizar o registro e a triagem taxonômica de comunidades ecológicas marinhas e estuarinas. Quanto aos objetivos, é exploratória -- ao buscar maior familiaridade com o problema por meio de entrevistas -- e descritiva -- ao comparar o desempenho dos métodos de triagem e identificação antes e depois da intervenção. Quanto aos procedimentos, trata-se de um estudo de campo com métodos mistos, que combina a análise qualitativa das entrevistas semiestruturadas com medições quantitativas de desempenho e de carga de trabalho mental, realizadas em pré-teste e pós-teste, conduzidas sob o referencial da IHC.

Operacionalmente, a pesquisa segue o processo iterativo do DT, cuja estrutura -- empatia, definição, ideação, prototipação e teste -- organiza as oito etapas descritas a seguir.

Na etapa de empatia, serão realizadas entrevistas semiestruturadas com pesquisadores de grupos de pesquisa que atuam no registro de comunidades ecológicas marinhas e estuarinas, selecionados por amostragem de conveniência. A participação será voluntária, mediante assinatura de Termo de Consentimento Livre e Esclarecido, com garantia de anonimato na análise dos relatos. O objetivo é compreender o fluxo de trabalho atual de coleta, triagem e identificação de táxons, abordando tanto o manuseio das amostras quanto a inclusão dos dados em planilhas até a análise desses dados. As entrevistas serão gravadas e transcritas, e as transcrições serão submetidas a análise temática: após a familiarização com os dados, gera-se um conjunto de códigos, a partir dos quais se constroem, revisam e definem os temas centrais. O produto desta etapa é um relatório de contexto organizado por temas, que mapeia as tarefas, as ferramentas e as dificuldades do método manual.

Em seguida, na etapa de definição, a análise temática conduzida anteriormente orienta a criação de personas e a construção de um mapa de empatia, técnicas de representação do usuário apresentadas na introdução. A partir dessas evidências, formula-se o POV. O produto desta etapa é o conjunto de artefatos de definição -- personas, mapa de empatia e POV, que direciona a ideação de soluções.

Antes do desenvolvimento, registra-se a linha de base (pré-teste): o desempenho com o método manual atual é medido a partir da execução, pelos mesmos pesquisadores, das mesmas tarefas de triagem e identificação de táxons, avaliando-se o tempo despendido por tarefa, a taxa de acerto na identificação e a carga de trabalho mental por meio do NASA-TLX. O produto desta etapa é o conjunto de dados de linha de base, que servirá de referência para a comparação posterior.

Na etapa de ideação, com base nos artefatos de definição, geram-se alternativas de solução para o registro e a triagem taxonômica, avaliadas junto aos participantes quanto à aderência ao fluxo de trabalho observado. O produto desta etapa é a alternativa de solução selecionada, com o escopo funcional inicial do sistema.

A partir da alternativa selecionada, procede-se à definição de requisitos, consolidando os requisitos funcionais e não funcionais do sistema mediante técnicas de engenharia de requisitos. O produto desta etapa é o documento de requisitos, que delimita o escopo da funcionalidade de triagem e análise de ocorrências taxonômicas a ser desenvolvida.

Segue-se a prototipação, na qual a alternativa é desenvolvida de forma incremental, com protótipos parciais testados junto aos participantes e ajustes incorporados conforme identificados. O produto desta etapa é a versão estável do sistema, pronta para a avaliação formal.

No pós-teste, os mesmos pesquisadores executam novamente as mesmas tarefas de triagem e identificação de táxons, agora utilizando o sistema desenvolvido. Medem-se a eficácia e a eficiência do uso, conforme definido pela ISO 9241-11, e a carga de trabalho mental por meio do NASA-TLX, em condições equivalentes às da linha de base. O produto desta etapa é o conjunto de dados de pós-teste, comparável ponto a ponto com os dados de pré-teste.

Por fim, na etapa de análise dos resultados, os dados obtidos antes e depois da intervenção são comparados por meio de estatística descritiva e de testes de comparação de amostras pareadas -- teste \textit{t} de Student pareado ou teste de Wilcoxon, conforme a distribuição verificada por teste de normalidade --, de modo a avaliar a significância estatística das diferenças observadas no tempo de tarefa, na taxa de acerto e na carga de trabalho mental. O produto desta etapa é a análise que confronta os resultados com os objetivos específicos da pesquisa.


\section{CRONOGRAMA}
O cronograma de atividades está apresentado no Quadro \ref{quadro_cronograma}, distribuído ao longo dos quatro semestres previstos para a conclusão do mestrado, entre os períodos de 2027.1 e 2028.2, e organizado segundo as etapas descritas na metodologia.

\begin{quadro}
\caption{\label{quadro_cronograma}Cronograma de atividades}
\begin{tabular}{|p{7cm}|c|c|c|c|}
	\hline
  \textbf{Atividade} & \textbf{2027.1} & \textbf{2027.2} & \textbf{2028.1} & \textbf{2028.2} \\ \hline
	Revisão bibliográfica & X & X & X &  \\ \hline
	Entrevistas semiestruturadas (empatia) & X & & & \\ \hline
	Construção de personas, mapa de empatia e POV (definição) & X & & & \\ \hline
	Coleta da linha de base (pré-teste) & X & & & \\ \hline
	Ideação de alternativas de solução & X & X & & \\ \hline
	Levantamento de requisitos & & X & & \\ \hline
	Prototipação incremental do sistema & & X & X & \\ \hline
	Coleta de dados de pós-teste & & & X & \\ \hline
	Análise estatística dos resultados & & & X & X \\ \hline
  Submissão de artigo científico & & & & X \\ \hline
	Defesa da dissertação & & & & X \\ \hline
\end{tabular}
\fonte{Autor.}
\end{quadro}

\cleardoublepage

\bibliography{plano_trabalho}
\end{document}
